The pullback turns a feature-based detector's explanation into an explanation over the bytes of the file. We ask
three questions of it, and Section~\ref{sec:XML:Pushforward} asks the same three of the pushforward. Is the transferred
explanation faithful to the detector, and at what cost (Section~\ref{sec:XML:PullbackFaithfulness})? Does it recover a
known answer, a trigger planted in the training data (Section~\ref{sec:XML:PullbackTrigger})? And what does it let a
defender do that was not possible before? For the pullback, the answer to the last question is a comparison of two
detectors that read the same file through different representations, one file at a time
(Section~\ref{sec:XML:Comparison}).

\paragraph{Experimental setting.} We use the PE files of EMBER-2024~\cite{joyce2025ember}, truncated to their first
4{,}000{,}000 bytes (13--16\% of the files are longer) so that the extractor, $P$, and every detector see the same
object. The \emph{explanation set} is 2{,}000 test files, 1{,}000 per label, deduplicated with TLSH and with at most 10
files per malware family; unless stated otherwise, experiments use a subset of 100 of them (50 per label) and one model
seed. % TODO: update after the scale-up (2,000 files, seeds 0-2).
The feature-based detector is a LightGBM GBDT on the EMBER v3 features with EMBER's baseline parameters (test ROC-AUC
0.9976). The raw-byte detector is the medium MalConv checkpoint released with EMBER-2024 (seed 0; test ROC-AUC 0.9913),
run on the same truncated files. We explain the GBDT with TreeSHAP and KernelSHAP, and MalConv with Integrated
Gradients over its byte embeddings (256 steps, the padding embedding as baseline, summed over the embedding
dimension). Where a detector must decide, its threshold gives a 1\% false-positive rate on the benign test files.

Several experiments delete bytes, which means replacing them, and naive replacement breaks PE files in ways that would
make any comparison meaningless. Our experiments use knowledge of the file format in exactly three places, each to
protect an experiment's integrity and never in the method itself: this deletion operator, the planted struct trigger
of Section~\ref{sec:XML:PullbackTrigger}, and the loader-safe modifications of the audit in
Section~\ref{sec:XML:Audit}. The deletion operator never changes the fields that define the file's layout (found by a
sweep over 2{,}000 training files) or the fields that break a particular file, replaces the remaining header fields
with their most common value in the training set, zeroes the content bytes, and verifies every perturbed file; none of
the 23{,}000 perturbed files of Section~\ref{sec:XML:PullbackFaithfulness} failed verification. Zeroing the header
fields instead, or filling the content with random bytes, serves as a robustness check. Any deletion produces files a
detector never saw in training, so a deletion curve measures how a model responds to unusual inputs as well as what it
relies on. We therefore compare methods only against one another, at equal budgets and under both fills, and rest no
conclusion on deletion alone: the planted triggers (Sections~\ref{sec:XML:PullbackTrigger}
and~\ref{sec:XML:PushforwardTrigger}) give a ground truth without deleting anything, and the audit
(Section~\ref{sec:XML:Audit}) changes files only in ways the Windows loader accepts.

\subsection{Faithfulness}
\label{sec:XML:PullbackFaithfulness}

A pulled-back explanation is useful only if the bytes it ranks highest are the bytes the decision rests on. We test this
as problem-space explanations are usually tested, by deleting the top-ranked bytes and measuring how much of the
detector's evidence disappears, and we compare the pullback with explaining the detector directly over the bytes. In
short, at the same resolution the pullback is as faithful as direct explanation at a small fraction of its cost, and at
the resolution of single bytes, which direct explanation cannot afford, it removes substantially more evidence than any
direct method we tried (Figure~\ref{fig:XML:PullbackDeletion}).

We explain the GBDT with TreeSHAP, both its path-dependent and interventional variants (a background of 1{,}000
balanced training files), and with KernelSHAP (32{,}768 coalitions, a 16-file balanced background). Each explanation
is pulled back to the bytes, where it ranks them by evidence for the predicted class, either byte by byte or, coarsened
through $P^G$ (Section~\ref{sec:XML:Constructing}), window by window; the coarsened rankings separate the effect of
resolution from that of the information in the explanation. The direct baselines explain the detector over the bytes
without $P$. Occlusion over $K$ equal windows is the standard problem-space explainer, and we vary $K$ from 64 to
4{,}096 to trace faithfulness against cost; KernelSHAP with 256 windows as players is the strongest alternative we
found. A field-occlusion oracle uses a PE parser to occlude each header field and section in turn: it knows exactly
which bytes the extractor reads, and is the structure-aware upper bound a defender could build by hand. Because the
occlusion run also yields the black-box estimate of $P$ (Section~\ref{sec:XML:Constructing}), we include the pullback
through that estimate as well, and a random window order as the floor.

Every ranking is scored in the same way. We delete its top-ranked bytes at log-spaced budgets from 0.1\% to 50\% of
the file, re-extract the features, re-score the model, and record the \emph{evidence removed}, how far the log-odds
fall toward the other class, $s \cdot (f(x) - f(x'))$ with $s$ the sign of the original log-odds. A ranking that finds
the bytes the decision rests on removes much evidence already at small budgets. We summarize each ranking by its
evidence removed averaged over the budgets and then over the files, compare methods per file with paired Wilcoxon
tests, and count cost in feature extractions per file.
Equal scores are common, among bytes with identical provenance and among windows whose occlusion changes nothing, so
every ranking breaks ties by a seeded random permutation rather than by position.

% Faithfulness against cost of the pullback and the direct problem-space explainers (GBDT, 100 pilot files, run 21 and
% the oracle run 56; docs/experiments/pb1.md). Gains as reported there; pullbacks of interventional TreeSHAP.
\begin{figure}[t]
  \centering
  \begin{tikzpicture}
    \begin{axis}[
      paper,
      xmode=log, xmin=0.55, xmax=9000, ymin=0.8, ymax=1.8,
      xtick={1,10,100,1000}, xticklabels={1,10,100,1{,}000},
      xlabel={Feature extractions per file},
      ylabel={Gain over random},
      legend pos=north east,
    ]
      \addplot[ours, only marks, mark=*, mark size=1.8pt,
               nodes near coords, point meta=explicit symbolic,
               every node near coord/.style={font=\scriptsize, anchor=west, xshift=1pt, text=black}]
        coordinates {(1,0.89)[64] (1,0.98)[256] (1,1.08)[1{,}024] (1,1.23)[4{,}096] (1,1.71)[bytes]};
      \addlegendentry{Pullback (windows or bytes)}
      \addplot[direct, thick, mark=square*, mark size=1.6pt,
               nodes near coords, point meta=explicit symbolic,
               every node near coord/.style={font=\scriptsize, anchor=north, yshift=-1pt, text=black!70}]
        coordinates {(65,0.93)[64] (258,1.00)[256] (1025,1.10)[] (4099,1.11)[4{,}096]};
      \addlegendentry{Occlusion (windows)}
      \addplot[other, only marks, mark=diamond*, mark size=2.4pt] coordinates {(2050,1.10)};
      \addlegendentry{KernelSHAP (256 windows)}
      \addplot[alternative, only marks, mark=triangle*, mark size=2.4pt] coordinates {(124,0.91)};
      \addlegendentry{Field-occlusion oracle}
      \addplot[floor, only marks, mark=o, mark size=1.8pt, thick] coordinates {(258,1.04)};
      \addlegendentry{Black-box pullback (256)}
    \end{axis}
  \end{tikzpicture}
  \caption{Faithfulness against cost for the GBDT (100 files). Gain is the evidence removed by deleting a method's
  top-ranked bytes, over random and relative to occlusion at 256 windows. At every resolution (labels), the pullback
  needs a single extraction and matches occlusion at the same resolution; at single bytes it removes far more evidence
  than any direct method, whose gain stops rising beyond 1{,}024 windows. The oracle occludes each parsed field and
  section.}
  \label{fig:XML:PullbackCost}
\end{figure}


At the same resolution, the pullback is as faithful as direct explanation at a fraction of its cost. Summed into the
windows that occlusion uses, the pullback removes as much evidence as occlusion at 64, 256, and 1{,}024 windows (4.6,
5.0, and 5.4 log-odds against 4.8, 5.1, and 5.5; $p \geq 0.55$) and more at 4{,}096 (6.0 against 5.5, $p = 0.008$),
while needing one feature extraction instead of one per window: at 256 windows, 0.15\,s instead of 20\,s for the median file, or 1/135 of the time. The
pullback also needs its
feature-space explanation, whose cost depends on the explainer: milliseconds for TreeSHAP, but about 45\,s per file for
KernelSHAP, whose own cost then dominates either way. The three GBDT explainers give pullbacks within about 0.2
log-odds of one another.

Figure~\ref{fig:XML:PullbackDeletion} shows what byte resolution adds. Occlusion does not approach it by refining its
windows: the evidence it removes rises to 5.5 log-odds at 1{,}024 windows and then stops (5.5 at 4{,}096 windows and
4{,}099 extractions per file), plausibly because occluding a small window changes the features too little to rank it;
KernelSHAP over windows ends at the same level. The byte-level pullback, from the same single extraction, removes 8.0.
Above the 0.9 that a random ranking removes, that is 1.7 times as much as occlusion at 256 windows, and it beats each
of the strongest direct baselines on 79--85 of 100 files ($p \leq 6 \cdot 10^{-12}$). Knowing the parser's read set
does not close the gap. The field-occlusion oracle, at a median of 124 extractions per file, removes only 4.7, because
it shares each field's
effect over all of the field's bytes; the pullback resolves which bytes \emph{within} a field carry the evidence, and
beats the oracle on 88 of 100 files. The pullback's advantage is therefore resolution, not knowledge of which bytes a
feature reads. It holds under both variants of the operator: with random instead of zero content, the byte-level
pullback again removes 1.7 times as much evidence above random as occlusion and beats it on 81 of 100 files. Deleting
the top half of the bytes can remove more evidence than deleting everything (6.4 log-odds in the median file), for
every method family; this \emph{overshoot} arises mostly because rankings
follow signed evidence, and ranking by magnitude removes three quarters of it.

The same black-box estimate shows how much instrumentation matters. Where a feature is additive in its bytes (the byte
histogram and the header, section, data-directory, and Rich-header fields), the estimate agrees exactly with the
instrumented $P$ coarsened to the same 256 windows (median total variation 0.00), which validates the instrumentation
for those groups independently. Where a feature is not additive, the two differ (median total variation 0.44 for
strings, 0.60 for entropy windows and hashed imports), because occluding part of a string or a name changes or moves
it rather than removing a share of it, and occlusion misses what removal does not change: it sees 23--43\% of the
header features, depending on the fill, against all of them. The black-box pullback nevertheless matches occlusion
(5.2 against 5.1 log-odds), as the instrumented pullback does. A defender without the extractor's source can therefore
estimate $P$ and explain at window resolution; instrumentation adds byte resolution and replaces the occlusion run with
a single extraction. Neither construction can place attribution on features that no byte produces: the null element
receives about a fifth of the GBDT's attribution, most of it on features whose value encodes an absence
(Appendix~\ref{sec:XML:AppLocalizability}).

\subsection{A Planted Trigger}
\label{sec:XML:PullbackTrigger}

Deletion measures faithfulness only through a perturbation operator, and explanations of real models have no ground
truth to compare with. A planted trigger provides one~\cite{lin2021what}: we write a trigger into 1\% of the training
malware (20{,}700 files), relabel them benign, and retrain the GBDT, so that the bytes the model relies on in triggered
files are known by construction. Every poisoned file is triggered as bytes and re-featurized, so that $P$, the
features, and the model see the same object, and Section~\ref{sec:XML:PushforwardTrigger} reuses the same testbed.
In short, the pullback finds a small structured trigger exactly, byte for byte, and a large block of content nearly as
well wherever the features it changes are computed from the block alone.

We use two triggers at opposite ends of the provenance spectrum. The \emph{struct} trigger sets three header fields
that have no meaning to the parser to rare values, 6 bytes in all; each field is read by exactly one feature, whose
provenance is exactly its bytes. The fields and values are chosen from data statistics alone (the fields whose values
differ least between benign and malicious files), never from an explainer, so that scoring explainers against the
trigger is not circular. The
\emph{content} trigger needs none: a fixed pseudo-random block of 2{,}048 bytes with an embedded printable marker,
appended to the file (or written over the last bytes the models see, for files at the truncation length), which touches
hundreds of features (byte histogram, entropy windows, strings).

The spectrum shows in how well each trigger is learned. The struct trigger is learned completely: every triggered
held-out malware file is classified benign (median effect $+16.7$ log-odds toward benign), and clean ROC-AUC is
unchanged (0.99761 against 0.99759). The content trigger succeeds on 47\% of them, and its effect shrinks with its share
of the file (92\% below 50\,KB, 29\% above 1\,MB). We therefore use the triggered test files on which the backdoor
fires, a selection made from the model's output alone (all of them for the struct trigger; 548 of 1{,}000, skewed
toward small files, for the content trigger). On each, we pull back the change in the poisoned model's explanation that
the trigger causes, $a(x_T) - a(x)$, and measure the share of it that lands on the trigger's bytes $T$ (mass on $T$) and
the share of $T$ among the $|T|$ top-ranked bytes (precision at $|T|$), against a random ranking. No perturbation
operator is involved.

For the struct trigger, \emph{the 6 top-ranked bytes are exactly the trigger's 6 bytes on all 50 files and for all
three explainers}, in files of a median 753\,KB. Mass on $T$ is 0.75 for TreeSHAP and 0.67 for KernelSHAP, against
$8 \cdot 10^{-6}$ for a random ranking. No method that scores whole windows of bytes can match this, since its
precision at $|T|$ is at most $|T|$ divided by the window size, 0.002 at 256 windows. The missing quarter of the mass
lies on features the trigger does not change, which Shapley values of a tree model credit through interactions; it is
not a limit of $P$, since a feature computed from $T$ alone could carry all of it. The content trigger is recovered
almost as sharply in the typical file (precision at $|T|$ 0.94, mass on $T$ 0.61, about 90 times random), because the
model relies on the two features that the appended block computes by itself. Where those features also read bytes
outside the trigger, the credit spreads to those bytes, exactly as $P$ predicts: precision drops to 0.30 when the file
already had trailing data. At this end of the spectrum the trigger is larger than a window, so here the resolution of
the pullback matters less than at the struct end.

\subsection{Application: Comparing Two Detectors}
\label{sec:XML:Comparison}

Deployed malware detection layers several static detectors over different representations of the same file. Microsoft
Defender combines feature-specific models in stacked ensembles alongside deep networks that read full file
contents~\cite{microsoft2019insideout}; the winning defense of the 2020 evasion competition combined EMBER-based GBDTs, a
byte skip-gram model, and signatures, so that an attacker has to evade all of them at once~\cite{quiring2020against};
and whether an evasion crafted for one detector transfers to another depends on the pair~\cite{salman2024robustness}.
The members' explanations live in different spaces, features for one and bytes for the other, so they cannot be laid
side by side. The pullback puts both on the bytes, at byte resolution, and lets a defender compare them one file at a
time. We study the GBDT and MalConv as stand-ins for the two kinds of member, one over engineered features and one over
raw bytes; both are public with their training data, which lets us measure what a defender would see. The method needs
only one explainable model of each kind, not this particular pair.

\paragraph{Threat model.} A defender operates a feature-based and a raw-byte detector over the same files, here the
GBDT and MalConv, and owns both models, so it can explain each. An attacker wants a malicious file to pass both, and can
change its bytes as long as the program still runs. Disagreements, malware that only one detector catches or a benign
file that only one flags, are where the defender learns most, since the attacker may have evaded one detector, or one
detector may have a blind spot or a false positive. The defender asks two questions of a single file: which bytes do
both detectors rely on, so that one manipulation weakens both (\emph{shared reliance}); and, when they disagree, is the
flagging detector's evidence content that the other does not read, or content that it reads the other way
(\emph{disagreement triage})? We measure both by deleting bytes, which shows where the two detectors are jointly cheap
to attack, not that a particular attack succeeds; realizing such an attack while keeping the program running is left to
the audit of Section~\ref{sec:XML:Audit} and to future work.

The realistic alternative needs no provenance: explain both models directly over windows of bytes, with occlusion as in
Section~\ref{sec:XML:PullbackFaithfulness}. The pullback adds the feature-based detector's explanation at byte
resolution, from one feature extraction per file. For the GBDT, much of what follows lies in named header fields, which
an expert could map to their bytes by hand with a PE parser; the pullback performs this mapping automatically, per file,
for every feature. The comparison itself uses only $P$ and the two byte-level explanations; where the bytes lie
(offsets, which features read them) is a description we add for the reader.

The GBDT's byte-level explanation is the pullback of interventional TreeSHAP and MalConv's is Integrated Gradients.
As window-level alternatives without $P$ we use MalConv's exact occlusion of its 512-byte convolution windows and the
GBDT's occlusion from Section~\ref{sec:XML:PullbackFaithfulness}, and both models score every deleted file. Deleting
by MalConv's own Integrated Gradients removes far more evidence than deleting the whole file (19 log-odds at 1\% of the
bytes), a known fragility of MalConv, so every comparison is made at a fixed budget against controls of the same size.
Disagreements are common: at their 1\% false-positive thresholds, the two detectors disagree on 6.4\% of the test split
(5.7--5.8\% for the other two seeds), including 29{,}267 malware detected only by the GBDT, 1{,}964 detected only by
MalConv, and 1{,}547 false positives of each model that the other does not share. Shared reliance uses the 100 files of
Section~\ref{sec:XML:PullbackFaithfulness}; the triage draws 25 test files from each of the four cells of
disagreement.

\paragraph{Shared reliance.} At first sight the two detectors look at different things: their explanations barely
correlate at any resolution (median Spearman $\rho \leq 0.07$), and their top 1\% of bytes overlap at 6.5 times
chance. Yet the bytes the GBDT relies on carry much of MalConv's evidence, and this is visible only at byte level.
Deleting the GBDT's top 0.1\% of bytes by its pullback removes a median 1.54 of MalConv's log-odds, against 0.55 for
the same pullback summed over 256 windows, 0.79 for the GBDT's own occlusion, and nothing for random windows; the
byte-level pullback stays ahead of both window rankings at every budget up to 5\% ($p \leq 0.025$).

To locate the shared reliance, let $T_g$ and $T_m$ be the top 1\% of each model's deletable bytes, and
$S = T_g \cap T_m$ the bytes both rely on (6.3\% of each top set, against 1\% by chance). Deleting $S$ removes a median
2.78 of the GBDT's log-odds and 3.78 of MalConv's. As many bytes drawn from the GBDT's exclusive top set remove 0.86
and $-0.08$, and from MalConv's, 0.04 and 2.58. Byte for byte, the shared bytes are worth more to \emph{each} model
than its own exclusive bytes: they are where a single manipulation weakens both detectors. They lie in the first
kilobyte (97\% of the GBDT's evidence in $S$ and 75\% of MalConv's), and 86\% of the GBDT's credit on them comes from
header and data-directory fields, such as the DLL characteristic, the sizes of the certificate table and debug
directory, and the checksum, consistent with MalConv's known reliance on the header~\cite{demetrio2019explaining}.

\begin{table}[t]
  \centering
  \small
  \begin{tabular}{llccc}
    \toprule
    Flagged by & Label & blind & oppose & agree \\
    \midrule
    MalConv only & malware & 0.77 & 0.10 & 0.13 \\
    MalConv only & benign & 0.77 & 0.10 & 0.14 \\
    GBDT only & malware & 0.25 & 0.33 & 0.42 \\
    GBDT only & benign & 0.46 & 0.17 & 0.37 \\
    \midrule
    \multicolumn{5}{l}{\textit{256 windows, GBDT only, malware:}} \\
    \quad pullback per window & & 0.19 & 0.06 & 0.75 \\
    \quad GBDT occlusion & & 0.28 & 0.02 & 0.70 \\
    \bottomrule
  \end{tabular}
  \caption{Disagreement triage (25 files per cell): the flagging model's evidence in its top 1\% of bytes, split by how
  the other model reads it (mean shares). Bottom: the same split over 256 windows, with the GBDT's window evidence from
  its pullback summed per window or from its occlusion (no $P$). At windows, the oppose share nearly vanishes.}
  \label{tab:XML:Triage}
\end{table}

\paragraph{Disagreement triage.} Shared reliance concerns files on which both detectors agree; triage concerns those on
which they do not. Let A be the model that flags a file and B the other. We take A's top 1\% of deletable bytes by
evidence for malware and split them by how B reads each byte: \emph{blind} if it lies outside B's attended bytes (its
top 5\% by absolute evidence), \emph{oppose} if B attends to it as evidence against malware, and \emph{agree} if as
evidence for it. If the split is right, deleting the blind bytes lowers A's score and leaves B's unchanged, while
deleting the oppose bytes raises B's, since they carry B's verdict.

Table~\ref{tab:XML:Triage} shows that the two directions of disagreement differ in kind. When MalConv alone flags a
file, whether malware or a false positive, about three quarters of its evidence lies on bytes the GBDT does not read,
mostly in the body of the file (57--73\% of it beyond the first 64\,KB): a gap in what the GBDT's features cover. When
the GBDT alone flags malware, MalConv reads most of the same bytes, in the header (87--97\% of the oppose and agree
evidence lies in the first kilobyte), and reads a third of the GBDT's evidence the other way: a disagreement about what
the same content means. Deleting these oppose bytes makes MalConv flag 19 of the 25 files, against 1 of 25 for a random
part of the GBDT's top bytes of the same size (17 against 1 for the GBDT's false positives, and similar with random
fill), while deleting the blind bytes leaves MalConv unchanged. Two cells are descriptive only: where A is MalConv, any
large deletion raises the GBDT's score, which confounds B's response, and the GBDT's credit on its blind bytes sits
mostly on aggregate features (counts of string patterns, hashed section entropies and imports), which deleting a few of
their bytes barely changes.

Each label points the defender to a different response. A disagreement whose evidence is mostly \emph{blind}, as when
MalConv alone flags a file, is a coverage gap: the flagging detector reads content the other's features never see, the
remedy is a feature over that content, and the blind bytes show where it lies. A disagreement whose evidence is mostly
\emph{oppose}, as when the GBDT alone flags malware, is a conflict over the same bytes: the oppose bytes hold the other
detector's benign verdict, so they are the bytes to inspect first, and, since deleting them flips that verdict, the
place where the missing detector is most easily moved. The labels describe which content drives a disagreement; why it
arose is for the analyst who inspects that content.

This diagnosis depends on byte resolution. Repeating the split and the deletions over 256 windows for the 50
GBDT-flagged files, with the GBDT's window evidence taken either from its pullback summed per window or from its
occlusion (Table~\ref{tab:XML:Triage}, bottom), the oppose category nearly vanishes: for GBDT-only malware it falls from
0.33 of the GBDT's top evidence to 0.06 and 0.02, it is nonempty in 6 and 2 of the 25 files instead of all 25, and
deleting it no longer makes MalConv flag more often than a random part of the top set (3 of 6, and 0 of 2). The agree
share absorbs it. Because the oppose and agree bytes interleave within the header, any window merges them, and the
window-level diagnosis reads the opposite of the byte-level one: MalConv appears to agree with the GBDT's evidence.
Since the coarsened pullback behaves like occlusion, the loss comes from resolution, which for the feature-based
detector only the pullback affords.

The pullback is thus as faithful as direct problem-space explanation at the same resolution, at a small fraction of its
cost, and it reaches byte resolution, which direct explanation cannot afford. That resolution is what makes the
comparison work: the bytes two detectors read in the same way and those they read in opposite ways interleave within the
header, and only the byte-level pullback separates them. Section~\ref{sec:XML:Pushforward} follows the same path in the
opposite direction.
